\section{Prompt templates}
\label{sec:all-prompts}
We list the prompt templates for LLM judges.
Please refer to our github repository~\footnote{\url{https://github.com/lm-sys/FastChat/tree/main/fastchat/llm_judge}} for full details.

\begin{figure}[h]
\centering









\includegraphics[width=\linewidth]{figures/default_prompt.pdf}
\caption{The default prompt for pairwise comparison.}
\label{fig:default-prompt-pair}
\end{figure}


\begin{figure}[h]
\centering






\includegraphics[width=\linewidth]{figures/single_prompt.pdf}
\caption{The default prompt for single answer grading.}
\label{fig:default-prompt-single}
\end{figure}


\begin{figure}[h]
\centering










\includegraphics[width=\linewidth]{figures/cot_prompt.pdf}
\caption{The chain-of-thought prompt for math and reasoning questions.}
\label{fig:cot-prompt-math}
\end{figure}

\begin{figure}[h]
\centering














\includegraphics[width=\linewidth]{figures/reference_prompt.pdf}
\caption{The prompt for reference-guided pairwise comparison.}
\label{fig:default-prompt-math}
\end{figure}

\begin{figure}[h]
\centering















\includegraphics[width=\linewidth]{figures/mt_prompt.pdf}
\caption{The prompt for multi-turn pairwise comparison.}
\label{fig:multiturn-prompt-pair}
\end{figure}

\begin{figure}[h]
\centering
\includegraphics[width=\linewidth]{figures/reference_mt_single_prompt.pdf}
\caption{The prompt for reference-guided multi-turn single-answer grading.}
\label{fig:reference-mt-single}
\end{figure}

\clearpage
\newpage
\section{Case Study}
We list several case studies. The examples are generated by \texttt{gpt-4-0314}. They may not be fully reproducible with future GPT-4 versions.

\begin{figure}[h]
\centering
\includegraphics[width=\linewidth]{figures/position_bias.pdf}
\caption{An example of position bias. When Assistant A is placed in the first position, GPT-4 thinks A is better, but its verdict changes when we swap the position of A and B. We observe similar pattern from other LLM judges such as Claude/GPT-3.5.}
\label{fig:position-bias-example}
\end{figure}


\begin{figure}[h]
\centering
\includegraphics[width=\linewidth]{figures/verbosity_bias.pdf}
\caption{An example of ``repetitive list'' attack to examine verbosity bias. Except for the two rephrased items (highlighted in red), Assistant A's answer is exactly the same as Assistant B. Both GPT-3.5 and Claude-v1 show a verbosity bias towards the longer and repetitive answer. Only GPT-4 successfully detected this attack.}
\label{fig:verbosity-bias-example}
\end{figure}



\begin{figure}[h]
    \centering
    \includegraphics[width=\linewidth]{figures/math_failure.pdf}
    \caption{With a default prompt, GPT-4 shows limited capability in grading math questions. Despite being able to answer the question itself, its judgment was influenced by the given answers, leading to arithmetic mistakes highlighted in yellow.}
    \label{fig:math-grading}
\end{figure}


\begin{figure}[h]
    \centering    \includegraphics[width=\linewidth]{figures/reasoning.pdf}
    \caption{An example of GPT-4's limited capability in grading reasoning question. Despite GPT-4 knows how to solve the question (if asked separately), it made a wrong judgement saying both assistants' wrong answers are correct.}
    \label{fig:reasoning-grading}
\end{figure}


\begin{figure}[h]
    \centering
    \includegraphics[width=\linewidth]{figures/cot_failure.pdf}
    \caption{An example of GPT-4's wrong judgment with chain-of-thought prompt. We can see GPT-4  exactly copied Assistant B's answer (which contains arithmetic errors) and determined Assistant A's answer is incorrect. This suggest GPT-4's chain-of-thought process can be significantly influenced by the given answers despite we ask it to think independently.}
    \label{fig:cot-failure}
\end{figure}


\begin{figure}[h]
    \centering
    \includegraphics[width=\linewidth]{figures/multiturn-failure.pdf}
    \caption{In this example, despite Assistant A correctly followed user's instruction to generate a concrete plan for the second example of its previous response, GPT-4 wrongly referred to  the second example in Assistant B's response, resulting in a wrong judgment. This suggests the prompt design that breaks the questions into two prompts may cause LLM judge struggle to locate assistants' previous responses.}
    \label{fig:mt-sep}
\end{figure}

\clearpage
\newpage
\section{Data Collection}
\label{sec:data-collection}
We describe our data collection process for both MT-bench and Chatbot Arena.

\subsection{MT-bench human evaluation}
\label{subsec:mt-data-collection}
\Cref{fig:screenshot-mt-bench} shows the normal voting interface. \Cref{fig:screenshot-mt-bench-agree} shows that we additionally show GPT-4's judgment to users and ask if it is reasonable when a human differs from GPT-4.

\begin{figure}[h]
    \centering
    \includegraphics[width=\linewidth]{figures/screenshot_mt_bench.png}
    \caption{The screenshot of MT-bench data collection. We show an instruction similar to the prompt we give to GPT-4. We present questions from MT-bench and answers from two random anonymous assistants and ask which one is better. We present the first-turn conversation and ask humans to vote, then repeat the same procedure for the second-turn. A user can skip up to 5 questions if they are not confident. For some questions (e.g., math, reasoning), they can also see a reference solution.}
    \label{fig:screenshot-mt-bench}
\end{figure}

\begin{figure}[H]
    \centering
    \includegraphics[width=\linewidth]{figures/screenshot_mt_bench_agree.png}
    \caption{The screenshot of MT-bench data collection. When human's vote differs from GPT-4, we additionally show GPT-4's judgment (red region in the screenshot) and ask the user to click one of the three buttons to decide whether GPT-4's judgment is reasonable.}
    \label{fig:screenshot-mt-bench-agree}
\end{figure}

To invite participants, we obtained their consent by letting them sign an application form. We pay them \$20 for judging 20 questions, which corresponds to an hourly rate of around \$35. The participants are mostly graduate students from more than ten universities.

\subsection{Chatbot Arena}
\label{subsec:chatbot-arena}
\Cref{fig:screenshot-arena} shows a screenshot of Chatbot Arena. Users are required to accept the terms of use, which obtain their consent and give us the right to release the conversation data.
The instructions are shown at the top of the interface. This is a free website. We do not pay users and any user can use this platform without registration. More introductions and analyses can be found at \url{https://lmsys.org/blog/2023-05-03-arena/}.

\begin{figure}[h]
    \centering
    \includegraphics[width=\linewidth]{figures/screenshot_arena.png}
    \caption{The screenshot of Chatbot Arena.}
    \label{fig:screenshot-arena}
\end{figure}

\subsection{Data Release}
We will clean the Personal Identifiable Information (PII) and tag toxic conversations with OpenAI moderation APIs for our dataset release.

\clearpage
\newpage
\section{Additional Experimental Results}
We present some additional experimental results.

\subsection{Position bias}
\label{subsec:additional-position-bias}
We test two more prompts and present the full results in \Cref{tab:position_bias_more_prompt}
``score'' changes the default prompt to let the model output two absolute scores instead of which one is better.
``short'' is a simplified version of our default prompt by removing instructions like ``Avoid any position bias..'', ``Begin your evaluation ... and provide a short explanation''.
We can find different prompts have different effects on different models.
For example, the "score" prompt can increase the consistency of GPT-3.5 but decreases it for Claude-v1 and GPT-4.

\begin{table}[h]
\centering
\footnotesize
\caption{Position bias on different models and prompts. Consistency is the percentage of cases where a judge gives consistent results when swapping the order of two assistants. ``Biased toward first'' is the percentage of cases when a judge favors the first answer. ``Error'' indicates wrong output formats. The two largest numbers in each column are in bold.}
\label{tab:position_bias_more_prompt}
\begin{tabular}{@{}llllll@{}}
\toprule
Judge & Prompt  & Consistency & Biased toward first & Biased toward second & Error \\
\midrule
\multirow{4}{*}{claude-v1} & default & 23.8\% & \textbf{75.0\%} & 0.0\% & 1.2\% \\
 & rename & 56.2\% & 11.2\% & \textbf{28.7\%} & \textbf{3.8\%} \\
 & score & 20.0\% & \textbf{80.0\%} & 0.0\% & 0.0\% \\
 & short & 22.5\% & \textbf{75.0\%} & 2.5\% & 0.0\% \\
\midrule
\multirow{4}{*}{gpt-3.5-turbo} & default & 46.2\% & 50.0\% & 1.2\% & 2.5\% \\
 & rename & 51.2\% & 38.8\% & 6.2\% & \textbf{3.8\%} \\
 & score & 55.0\% & 33.8\% & \textbf{11.2\%} & 0.0\% \\
 & short & 38.8\% & 57.5\% & 3.8\% & 0.0\% \\
\midrule
\multirow{4}{*}{gpt-4} & default & \textbf{65.0\%} & 30.0\% & 5.0\% & 0.0\% \\
 & rename & \textbf{66.2\%} & 28.7\% & 5.0\% & 0.0\% \\
 & score & 51.2\% & 46.2\% & 2.5\% & 0.0\% \\
 & short & 62.5\% & 35.0\% & 2.5\% & 0.0\% \\
\bottomrule
\end{tabular}
\end{table}

As shown in \Cref{tab:position_bias_category}, position bias is more noticeable on open questions like writing and stem/humanity knowledge questions.
On math and coding questions, LLM judges are more confident even though their judgments can often be wrong, as we show in \Cref{subsec:llmlimit}. Finally, we study how the model pairs influence position bias by using GPT-4 and the default prompt to judge three different model pairs. As shown in \Cref{tab:position_bias_pairs}, the position bias is more noticeable for models with close performance and can almost disappear when the performance of the two models differs a lot.

\begin{table}[h]
\centering
\footnotesize
\caption{Position bias on different categories. The two largest numbers in each column are in bold.}
\label{tab:position_bias_category}
\begin{tabular}{lllll}
\toprule
Category & Consistent & Biased toward first & Biased toward second \\
\midrule
writing & 42.0\% & 46.0\% & 12.0\% \\
roleplay & 68.0\% & 30.0\% & 2.0\% \\
reasoning & 76.0\% & 20.0\% & 4.0\% \\
math & \textbf{86.0\%} & 4.0\% & \textbf{10.0\%} \\
coding & \textbf{86.0\%} & 14.0\% & 0.0\% \\
extraction & 78.0\% & 12.0\% & \textbf{10.0\%} \\
stem & 44.0\% & \textbf{54.0\%} & 2.0\% \\
humanities & 36.0\% & \textbf{60.0\%} & 4.0\% \\
\bottomrule
\end{tabular}
\end{table}

\begin{table}[h]
\centering
\footnotesize
\caption{Position bias on different model pairs.}
\label{tab:position_bias_pairs}
\begin{tabular}{lllll}
\toprule
Pair & Consistent & Biased toward first & Biased toward second \\
\midrule
GPT-3.5 vs Claude-V1  & 67.5\% & 23.8\% & 8.8\% \\
GPT-3.5 vs Vicuna-13B & 73.8\% & 23.8\% & 2.5\% \\
GPT-3.5 vs LLaMA-13B  & \textbf{98.8\%} & 1.2\% & 0.0\% \\
\bottomrule
\end{tabular}
\end{table}

\subsection{Few-shot judge}
\label{subsec:additional-few-shot-judge}
We examine how few-shot examples improve LLM judges. As shown in \Cref{tab:few_shot_judge_consistency}, they improve the consistency of all three LLM judges significantly. It almost alleviates the position bias of GPT-4, but moves the position bias of GPT-3.5 from the first position to the second position. We then measure the agreement between few-shot GPT-4 pairwise comparison and humans on MT-bench, but found it performs similarly to zero-shot GPT-4 pairwise comparison.

\begin{table}[h]
\centering
\footnotesize
\caption{Improvements of the few-shot judge on consistency for position bias.}
\label{tab:few_shot_judge_consistency}

\begin{tabular}{@{}llllll@{}}
\toprule
Model & Prompt  & Consistency & Biased toward first & Biased toward second & Error \\
\midrule
\multirow{2}{*}{Claude-v1} & zero-shot & 23.8\% & 75.0\% & 0.0\% & 1.2\% \\
 & few-shot & \textbf{63.7\%} & 21.2\% & 11.2\% & 3.8\% \\
\midrule
\multirow{2}{*}{GPT-3.5} & zero-shot & 46.2\% & 50.0\% & 1.2\% & 2.5\% \\
 & few-shot & \textbf{55.0\%} & 16.2\% & 28.7\% & 0.0\% \\
\midrule
\multirow{2}{*}{GPT-4} & zero-shot & 65.0\% & 30.0\% & 5.0\% & 0.0\% \\
 & few-shot & \textbf{77.5\%} & 10.0\% & 12.5\% & 0.0\% \\
\bottomrule
\end{tabular}
\end{table}

\subsection{Agreement Evaluation}
\label{subsec:agreement-evaluation}

\textbf{Agreement calculation.}
We define the agreement between two types of judges as the probability of randomly selected individuals (but not identical) of each type agreeing on a randomly selected question.
For example, if we are comparing GPT-4 and Claude, the agreement is the probability of GPT-4 and Claude agreeing on the vote for a randomly selected question.
If we are comparing GPT-4 and humans, the agreement is the probability of GPT-4 and a randomly selected human agreeing on the vote for a randomly selected question.
The agreement among humans themselves is the probability of two randomly selected but not identical humans agreeing on the vote for a randomly selected question.

Note that the agreement among humans could be a lower estimation compared to the agreement of GPT4 and humans.
Consider three humans who voted ``A'', ``A'', and ``B'' for a question, respectively.
The agreement among them is only $\frac{1}{3}$, as there are three pairs ``(A, A)'', ``(A, B)'', and ``(A, B)''.
But the agreement between GPT4 and those three is $\frac{2}{3}$ if GPT4 voted ``first'' and $\frac{1}{3}$ otherwise.

Therefore, to have a more comprehensive understanding of what happened, we introduce a new judge type called human-majority, which considers the majority of human votes for each question.
The agreement between GPT4 and human-majority is then calculated as the probability of GPT4 agreeing with the majority of human votes on a randomly selected question.
\textit{The upper bound of the agreement between GPT-4 and humans is the agreement between human-majority and human.}
When there is no majority vote for a question, the agreement is counted by an even split.
For example, if there are an equal number of ``A'' and ``B'' human votes for a question, and GPT4 votes ``A'', the agreement is counted as $\frac{1}{2}$ on this question.

\textbf{More results.} \Cref{table:agreement_mt_full} shows more agreement results on MT-bench. In addition to expert labelers (denoted as ``Human''), we also include author votes (denoted as ``Author'').

\begin{table}[h]
\centering
\vspace{-1em}
\caption{Agreement between two types of judges on MT-bench. ``G4-P'' and ``G4-S'' denote GPT-4 with pairwise comparison and single-answer grading, respectively.
``C'' denotes Claude.
``Human'' denotes expert labelers (excluding authors). `Human-M'' denotes the majority vote of humans.
The single-answer grading can be converted into pairwise comparison results for calculating the agreement.
We report two setups: ``S1'' includes non-tie, tie, and inconsistent (due to position bias) votes and counts inconsistent as a tie; ``S2'' only includes non-tie votes.
The agreement between two random judges under each setup is denoted as ``R=''.
The top value in each cell is the agreement, and the bottom gray value is \#votes.
}
\footnotesize
\begin{subtable}[t]{\textwidth}
\resizebox{1\columnwidth}{!}{
\begin{tabular}[t]{lrrrrrrrrrr}
\toprule
Setup & \multicolumn{5}{c}{S1 (R = 33\%)} & \multicolumn{5}{c}{S2 (R = 50\%)}  \\
\cmidrule(lr){2-6}\cmidrule(lr){7-11}

Judge & G4-S & C & Author & Human & Human-M & G4-S & C & Author & Human & Human-M \\
\midrule

 G4-P & \shortstack{70\% \\ \textcolor{gray}{ 1138 }} & \shortstack{63\% \\ \textcolor{gray}{ 1198 }} & \shortstack{69\% \\ \textcolor{gray}{ 345 }} & \shortstack{66\% \\ \textcolor{gray}{ 1343 }} & \shortstack{67\% \\ \textcolor{gray}{ 821 }} & \shortstack{97\% \\ \textcolor{gray}{ 662 }} & \shortstack{94\% \\ \textcolor{gray}{ 582 }} & \shortstack{92\% \\ \textcolor{gray}{ 201 }} & \shortstack{85\% \\ \textcolor{gray}{ 859 }} & \shortstack{85\% \\ \textcolor{gray}{ 546 }}\\
 \midrule
G4-S & - & \shortstack{66\% \\ \textcolor{gray}{ 1136 }} & \shortstack{67\% \\ \textcolor{gray}{ 324 }} & \shortstack{60\% \\ \textcolor{gray}{ 1280 }} & \shortstack{60\% \\ \textcolor{gray}{ 781 }} & - & \shortstack{90\% \\ \textcolor{gray}{ 563 }} & \shortstack{94\% \\ \textcolor{gray}{ 175 }} & \shortstack{85\% \\ \textcolor{gray}{ 739 }} & \shortstack{85\% \\ \textcolor{gray}{ 473 }}\\
 \midrule
C & - & - & \shortstack{58\% \\ \textcolor{gray}{ 343 }} & \shortstack{54\% \\ \textcolor{gray}{ 1341 }} & \shortstack{55\% \\ \textcolor{gray}{ 820 }} & - & - & \shortstack{89\% \\ \textcolor{gray}{ 141 }} & \shortstack{85\% \\ \textcolor{gray}{ 648 }} & \shortstack{86\% \\ \textcolor{gray}{ 414 }}\\
 \midrule
Author & - & - & \shortstack{69\% \\ \textcolor{gray}{ 49 }} & \shortstack{65\% \\ \textcolor{gray}{ 428 }} & \shortstack{55\% \\ \textcolor{gray}{ 93 }} & - & - & \shortstack{87\% \\ \textcolor{gray}{ 31 }} & \shortstack{83\% \\ \textcolor{gray}{ 262 }} & \shortstack{76\% \\ \textcolor{gray}{ 46 }}\\
 \midrule
Human & - & - & - & \shortstack{63\% \\ \textcolor{gray}{ 721 }} & \shortstack{81\% \\ \textcolor{gray}{ 892 }} & - & - & - & \shortstack{81\% \\ \textcolor{gray}{ 479 }} & \shortstack{90\% \\ \textcolor{gray}{ 631 }}\\
 \midrule
\end{tabular}
}
\caption{First Turn}
\label{table:mt_bench_more_first_turn}
\end{subtable}

\begin{subtable}[t]{\textwidth}
\centering
\begin{tabular}[t]{lrrrrrrrr}
\toprule
Setup & \multicolumn{4}{c}{S1 (R = 33\%)} & \multicolumn{4}{c}{S2 (R = 50\%)} \\
\cmidrule(lr){2-5}\cmidrule(lr){6-9}

Judge & G4-S & Author & Human & Human-M & G4-S & Author & Human & Human-M \\
\midrule
 
G4-P & \shortstack{70\% \\ \textcolor{gray}{ 1161 }} & \shortstack{66\% \\ \textcolor{gray}{ 341 }} & \shortstack{66\% \\ \textcolor{gray}{ 1325 }} & \shortstack{68\% \\ \textcolor{gray}{ 812 }} & \shortstack{95\% \\ \textcolor{gray}{ 727 }} & \shortstack{88\% \\ \textcolor{gray}{ 205 }} & \shortstack{85\% \\ \textcolor{gray}{ 864 }} & \shortstack{85\% \\ \textcolor{gray}{ 557 }}\\
 \midrule
G4-S & - & \shortstack{65\% \\ \textcolor{gray}{ 331 }} & \shortstack{59\% \\ \textcolor{gray}{ 1285 }} & \shortstack{61\% \\ \textcolor{gray}{ 783 }} & - & \shortstack{89\% \\ \textcolor{gray}{ 193 }} & \shortstack{84\% \\ \textcolor{gray}{ 776 }} & \shortstack{85\% \\ \textcolor{gray}{ 506 }}\\
 \midrule
Author & - & \shortstack{67\% \\ \textcolor{gray}{ 49 }} & \shortstack{68\% \\ \textcolor{gray}{ 413 }} & \shortstack{63\% \\ \textcolor{gray}{ 87 }} & - & \shortstack{87\% \\ \textcolor{gray}{ 31 }} & \shortstack{86\% \\ \textcolor{gray}{ 273 }} & \shortstack{84\% \\ \textcolor{gray}{ 54 }}\\
 \midrule
Human & - & - & \shortstack{67\% \\ \textcolor{gray}{ 707 }} & \shortstack{83\% \\ \textcolor{gray}{ 877 }} & - & - & \shortstack{82\% \\ \textcolor{gray}{ 474 }} & \shortstack{91\% \\ \textcolor{gray}{ 629 }}\\
 \midrule
\end{tabular}
\caption{Second Turn}
\label{table:mt_bench_more_second_turn}
\end{subtable}
\label{table:agreement_mt_full}
\end{table}



 


\subsection{Category-wise scores with single-answer grading}
\label{subsec:category-wise-scores}
We use single-answer grading to evaluate 6 models on MT-bench and plot the category-wise scores in \Cref{fig:mt-bench-single-cat}.

\begin{figure}[t]
    \centering
    \includegraphics[width=0.8\linewidth]{figures/mtradar.pdf}
    \caption{Category-wise scores of 6 models on MT-bench.}
    \label{fig:mt-bench-single-cat}
\end{figure}

\clearpage
\newpage
\section{Training Details of Vicuna Models}
\label{sec:training-details}

Vicuna is created by fine-tuning a LLaMA base model using user-shared conversations gathered from ShareGPT.com with its public APIs.
ShareGPT is a website where users can share their ChatGPT conversations.
To ensure data quality, we convert the HTML back to markdown and filter out some inappropriate or low-quality samples, which results in 125K conversations after data cleaning.\footnote{In this study, we use more data (125K) than the version in our earlier blog post (70K).}
We then divide lengthy conversations into smaller segments that fit the model's maximum context length.

We construct three training datasets with different scales from this cleaned ShareGPT dataset. Their statistics are in \Cref{table:dataset_composition}, where we also compare it with Alpaca~\cite{alpaca} dataset. ``All'' is the full dataset. ``Single'' only includes the first turn of each conversation. ``Selected'' is a small high-quality dataset of 3K sequences. To construct the ``Selected'' dataset, we pick sequences that include at least 3 turns of conversations generated by GPT-4 and run a clustering algorithm to divide them into 3K clusters and pick the centroid of each cluster.

All models (Vicuna-7B/13B) are trained with the same hyperparameters: global batch size=128, learning=2e-5, epochs=3, seq length=2048. Except for ``Selected'', which we train for 5 epochs.
The training code is built on top of the Alpaca code but additionally handles multi-turn conversations.
The training is done with 8x A100 GPUs. The longest single training run takes around 2 days.
We utilize SkyPilot~\cite{skypilot} managed spot instances for saving training costs and FlashAttention~\cite{dao2022flashattention} for memory optimizations.
The training code is available at \url{https://github.com/lm-sys/FastChat}.


\begin{table}[h]
\centering
\caption{Dataset statistics}
\label{table:dataset_stats}
\footnotesize
\begin{tabular}{llllll}
\toprule
Dataset Name                 & Alpaca  & Selected & Single   & All   \\
\midrule
\#Token                      & 4.4M    & 4.8M     & 184M     & 370M  \\
\#Sequence                   & 52K     & 3K       & 257K     & 257K  \\
Avg. turns of conversation   & 1.0     & 4.0      & 1.0      & 2.9   \\
Avg. response length (token) & 65      & 343      & 473      & 373   \\
\bottomrule
\end{tabular}
\end{table}

\clearpage
\newpage
\section{Exploring Vicuna as a judge}
\label{sec:vicuna-judge}
In this paper, we mostly evaluate the ability of close-sourced models such as GPT-4 as a proxy for human evaluations. However, model services such as GPT-4 can also become expensive with a growing number of evaluations. On the other hand, popular open-sourced LLMs, e.g. Vicuna-13B shows strong language understanding capability, and are much cheaper than close-sourced LLMs. In this section, we further explore the potential of using Vicuna-13B as a more cost-friendly proxy. 

\subsection{Zero-Shot Vicuna}
When using as-it-is (zero-shot), Vicuna-13B noticeably suffers from limitations we discuss, e.g. position bias. As shown in \Cref{tab:position_bias_vicuna}, Vicuna-13B has a consistency rate from 11.2\% to 16.2\% across different prompt templates, much lower than all the closed-sourced models. In addition, it has a high error rate (from 22.5\% to 78.8\%) because of its weaker instruction-following capability. In many scenarios, Vicuna-13B provides responses such as "Answer A is better than answer B", without following the pre-defined template. These responses are rendered as natural languages and are difficult to be parsed automatically, making the model less useful in a scalable and automatic evaluation pipeline. 

\subsection{Arena Fine-tuned Vicuna}
\paragraph{Training}
Due to the incapability of the zero-shot Vicuna-13B model, we further finetune the model with human votes from Chatbot Arena. Specifically, we randomly sample 22K single-turn votes from the arena, covering all models supported by the time of this paper submission (GPT-4, GPT-3.5, Claude-v1, Vicuna-13b, Vicuna-7b, Koala-13B, Alpaca-13B,LLaMA-13B, Dolly-12B, FastChat-T5, RWKV-4-Raven, MPT-Chat, OpenAssistant, ChatGLM, and StableLM), to expose the model with a wider range of chatbot outputs and human preferences. We use 20K votes for training, and 2K for validation. To address the aforementioned weak instruction following problem, we formulate the problem as a 3-way sequence classification problem. Thus, the model simply needs to predict which one of the chat-bot outputs is better (or tie), without needing to exactly following the provided answer template.
In particular, we construct an input by using the default prompt and the two model answers. The labels are A, B, and tie (including both-bad-vote and tie-vote).
We train for 3 epochs with a cosine learning rate scheduler and a 2e-5 maximum learning rate. We use the 2K validation dataset to choose hyper-parameters, and test on the same 3K dataset in the main body of the paper.

\paragraph{Position bias results}
The results for position bias are provided in \Cref{tab:position_bias_vicuna}.
The consistency improves significantly from 16.2\% to 65.0\%. Due to the classification formulation, every output is recognizable (error rate 0\%).
In addition, we measure the classification accuracy over the test dataset.

\paragraph{Agreement results}
It achieves 56.8\% when including all three labels, and 85.5\% when excluding tie predictions and labels, significantly outperforming random guesses of 33\% and 50\% respectively, and show positive signals to match GPT-4 (66\% and 87\% respectively).
In conclusion, a further fine-tuned Vicuna-13B model shows strong potential to be used as a cheap open-sourced replacement for expensive closed-sourced LLMs. A similar conclusion is also found by a concurrent paper\cite{wang2023pandalm}. 

\begin{table}[h]
\centering
\footnotesize
\caption{Position bias of the Vicuna-13B model without and with further fine-tuning. We denote them as Vicuna-13B-Zero-Shot and Vicuna-13B-Fine-Tune respectively. Consistency is the percentage of cases where a judge gives consistent results when swapping the order of two assistants. ``Biased toward first'' is the percentage of cases when a judge favors the first answer. ``Error'' indicates wrong output formats. The largest number in each column is in bold.}
\label{tab:position_bias_vicuna}

\begin{tabular}{@{}llllll@{}}
\toprule
Judge & Prompt  & Consistency & Biased toward first & Biased toward second & Error \\
\midrule
\multirow{3}{*}{Vicuna-13B-Zero-Shot} & default & 15.0\% & \textbf{53.8\%} & 8.8\% & 22.5\% \\
 & rename & 16.2\% & 12.5\% & \textbf{40.0\%} & 31.2\% \\
 & score & 11.2\% & 10.0\% & 0.0\% & \textbf{78.8\%} \\ 
\midrule
\multirow{1}{*}{Vicuna-13B-Fine-Tune} & default & \textbf{65.0\%} & 27.5\% & 7.5\% & 0.0\% \\
\bottomrule
\end{tabular}
\end{table}

